\section{Derating via Below-Rated Pitch Control}

\subsection{Theory}

Derating (downregulation) reduces a turbine's power output by increasing the pitch
angle, moving the operating point to a lower effective wind speed on the turbine's
$P(u)$ and $C_T(u)$ curves.
Given a free-stream (or wake-effective) wind speed $u_\infty$ and a desired power
fraction $f \in [0,1]$, the operating wind speed $u_\text{op}$ is found by solving

\begin{equation}
    P(u_\text{op}) = f \cdot P(u_\infty), \qquad u_\text{op} \leq u_\infty,
    \label{eq:derate_inversion}
\end{equation}

where $P(\cdot)$ is the tabular power curve of the turbine.
The corresponding derated thrust coefficient and power output are then

\begin{equation}
    C_T^\text{d} = C_T(u_\text{op}), \qquad P^\text{d} = P(u_\text{op}).
    \label{eq:derated_ct}
\end{equation}

Because $u_\text{op} < u_\infty$, the turbine operates at a lower induction, producing a
weaker wake deficit downstream.
This is physically correct: derating follows the real aerodynamic relationship between
pitch angle, $C_T$, and power, rather than applying an arbitrary scaling to thrust.

Equation~\eqref{eq:derate_inversion} is solved by inverting the power curve on its
monotonically increasing portion (cut-in to rated wind speed).
Above rated, $P$ is flat so the inverse is not unique; however, $C_T$ is also
approximately constant in that region, making the distinction unimportant for wake
modelling purposes.

The derate fraction $\delta \in [0,1]$ relates to $f$ by

\begin{equation}
    f = 1 - \delta,
\end{equation}

so $\delta = 0$ means no curtailment ($f=1$, full power) and $\delta = 1$ means
full shutdown ($f=0$, turbine off).

\subsection{Effect on Downstream Wakes}

In the Dynamic Wake Meandering (DWM) model, the wake deficit and added turbulence
intensity downstream of turbine $i$ scale with $C_T^{(i)}$.
Derating reduces $C_T^{(i)}$ via~\eqref{eq:derated_ct}, which shrinks the DWM wake
particles emitted at turbine $i$, resulting in a higher effective wind speed at
downstream rotors.
The net farm power change is

\begin{equation}
    \Delta P_\text{farm} = \underbrace{-\delta \cdot P(u_\infty^{(i)})}_{\text{upstream loss}}
    + \underbrace{\sum_{j \in \text{wake}(i)} \bigl[P(u_j^\text{d}) - P(u_j)\bigr]}_{\text{downstream gain}},
\end{equation}

where $u_j^\text{d} > u_j$ is the improved wind speed at turbine $j$ under derating.
Derating is only beneficial when the downstream gain exceeds the upstream loss.

\subsection{Implementation in WindGym}

\paragraph{Power-curve inversion (\texttt{DeratingModel}).}
A new \texttt{AdditionalModel} subclass (\texttt{WindGym/core/derating.py}) is appended
to the turbine's \texttt{PowerCtTabular.model\_lst}.
At construction time the monotonically increasing portion of the power curve
$\{u_k, P_k\}$ is stored.
At each call, given $u_\infty$ and $\delta$, equation~\eqref{eq:derate_inversion} is
solved via \texttt{np.interp} on the pre-inverted table, and $C_T$ and $P$ are
evaluated at $u_\text{op}$.
The model is attached once with \texttt{add\_derating\_to\_turbine(turbine)} and is
idempotent across environment resets.

\paragraph{Action space.}
When \texttt{derate\_action: true} in the environment config, the action space doubles
to $2n_\text{turb}$ entries:

\begin{equation}
    \mathbf{a} = [\underbrace{\psi_0,\ldots,\psi_{n-1}}_{\text{yaw offsets}},\;
                  \underbrace{a_0^\delta,\ldots,a_{n-1}^\delta}_{\text{derate actions}}],
    \quad a_i^\delta \in [-1,1].
\end{equation}

Each derate action is mapped to $\delta_i \in [\delta_\min, \delta_\max]$ by

\begin{equation}
    \delta_i = \frac{a_i^\delta + 1}{2} \cdot (\delta_\max - \delta_\min) + \delta_\min.
\end{equation}

\paragraph{Sensor hookup (Dynamiks backend).}
A \texttt{derate} sensor is registered on \texttt{PyWakeWindTurbines}.
At each environment step, \texttt{wts.sensors.derate} is set to the current
$[\delta_0,\ldots,\delta_{n-1}]$ vector.
Dynamiks' \texttt{get\_kwargs()} propagates the sensor values automatically to the
\texttt{powerCtFunction}, so the updated $C_T^{(i)}$ values are used when the DWM
particles are emitted in the next time step.

\paragraph{Wake particle coverage.}
For side-by-side layouts all turbines share the same downwind position, giving a
farm downwind extent of zero.
Dynamiks' default particle count (based on farm extent) therefore collapses to its
minimum of 10 particles, covering only $10 \times 0.2D \approx 2D$ downstream.
WindGym now ensures at least $15D$ of coverage by computing

\begin{equation}
    N_\text{particles} = \max\!\left(\left\lceil
        \frac{\max(1.2\,L_x,\;15D)}{d_\text{particle}\,D}
    \right\rceil,\; 10\right),
\end{equation}

where $L_x$ is the farm downwind extent and $d_\text{particle} = 0.2$ (normalised
particle spacing).

\paragraph{Reward.}
No special reward term is required.
Unnecessary derating directly reduces farm power, so the standard power reward
already penalises it.
The agent must discover autonomously that upstream derating is only worthwhile when
the downstream power recovery exceeds the upstream sacrifice.
